\section[Conclusão]{Conclusão}

Tanto por minha experiência limitada com as tecnologias abordadas quanto por
ser minha primeira AGES, este projeto foi incrivelmente intimidador e
desafiador. Minha timidez e meu receio de não conseguir realizar algumas tarefas
também contribuíram para essa dificuldade.

Certamente, se eu tivesse abordado esse problema no início do semestre,
poderia ter ajudado o time de forma mais eficaz. Acho que alguns colegas
enfrentaram a mesma timidez, pois havia muito silêncio entre os alunos AGES I
e alguns AGES II. Isso foi uma pena, ainda mais porque muitos colegas,
incluindo alunos AGES II, III e IV, sempre foram prestativos e comprometidos
com o projeto.

Além dos desafios internos enfrentados pela equipe, tivemos problemas já na
Sprint 0 ao tentar levantar os requisitos com as stakeholders e definir o escopo
do projeto.
Devido à falta de clareza sobre o propósito da primeira reunião, não conseguimos
levantar nenhum requisito nesse dia, resultando em um atraso de uma semana no
início do desenvolvimento. Adicionalmente, o escopo, que foi delineado após a
reunião, mostrou-se extremamente aberto e sujeito a mudanças durante novos
encontros com os clientes.

Embora eu tenha sido crítico com a disciplina de AGES I e este projeto tenha
se mostrado confuso e complexo, aprendi muito com meus colegas. Afinal,
superar desafios é essencial ao desenvolvimento,
pois os obstáculos também oferecem oportunidades de progredir.

Levo comigo a lição de que a comunicação é um dos aspectos mais relevantes do
trabalho em grupo. Embora tenham sido escolhidas tecnologias complexas e
avançadas, as dificuldades do projeto teriam sido resolvidas mais rapidamente
se a comunicação do time fosse melhor.

Por fim, noto que aprimorei minha capacidade de interagir e me expressar com
meus colegas no ambiente de trabalho. Tenho certeza de que esses aprendizados
serão importantes para minhas futuras experiências na AGES.
